\documentclass[11pt]{article}
\usepackage[a4paper,margin=22mm]{geometry}
\usepackage{amsmath,booktabs,longtable,hyperref}
\begin{document}
\title{Normalized amplitude features: 12 equations}\author{PCG-BP analysis version 2}\date{}\maketitle
For each complete adjacent S1-S2-S1 cycle, a, b, c, d, e denote S1 onset, S1 offset, S2 onset, S2 offset and next S1 onset, in zero-based recording samples. They must satisfy a < b < c < d < e. Sampling frequency is f\_s. States use half-open sample intervals [a,b), [b,c), [c,d), [d,e).

Only complete cycles contained in a segment contribute. The next S1 onset may equal the exclusive segment end; its waveform belongs to the next segment. Detection is performed once over the whole recording.

Normalize the filtered waveform separately within each saved segment to zero mean and unit population SD. L is the number of segment samples. A\_Q is the mean absolute normalized waveform amplitude over state Q, not an envelope peak or RMS. Exclude a cycle if any state amplitude is at most 1e-12. All six ratios use the same valid-cycle pool.
\[z[n]=\frac{x[n]-\mu_x}{\sigma_x},\quad \mu_x=\frac{1}{L}\sum_n x[n],\quad \sigma_x=\sqrt{\frac{1}{L}\sum_n(x[n]-\mu_x)^2}\]
\[A_{Q,i}=\frac{1}{|I_{Q,i}|}\sum_{n\in I_{Q,i}}|z[n]|\]
\begin{longtable}{lll}\toprule Quantity & Equation & Output columns\\\midrule
Amp\_SysS1 & $A_{Sys,i}/A_{S1,i}$ & m\_Amp\_SysS1, sd\_Amp\_SysS1\\
Amp\_DiaS2 & $A_{Dia,i}/A_{S2,i}$ & m\_Amp\_DiaS2, sd\_Amp\_DiaS2\\
Amp\_S1S2 & $A_{S1,i}/A_{S2,i}$ & m\_Amp\_S1S2, sd\_Amp\_S1S2\\
Amp\_S1Dia & $A_{S1,i}/A_{Dia,i}$ & m\_Amp\_S1Dia, sd\_Amp\_S1Dia\\
Amp\_SysDia & $A_{Sys,i}/A_{Dia,i}$ & m\_Amp\_SysDia, sd\_Amp\_SysDia\\
Amp\_S2Sys & $A_{S2,i}/A_{Sys,i}$ & m\_Amp\_S2Sys, sd\_Amp\_S2Sys\\
\bottomrule\end{longtable}
\section*{Aggregation within each segment}
\[m_q=\frac{1}{K}\sum_{i=1}^{K}q_i,\qquad sd_q=\sqrt{\frac{1}{K-1}\sum_{i=1}^{K}(q_i-m_q)^2}.\]
Ratios are calculated per cycle, then averaged. Sample SD uses $K-1$, an explicit implementation convention. Mean is NaN for $K=0$; SD is NaN for $K<2$. Invalid or ambiguous cycles are excluded, never bridged or filled with zero.

\section*{Reference and adaptations}
Tang et al. (2018), \emph{PCG Classification Using Multidomain Features and SVM Classifier}, Tables 1--2 and Sections 2.3--2.5.3. \url{https://doi.org/10.1155/2018/4205027}.
The feature definitions follow the paper. This project retains its 30--500 Hz Butterworth filter and timing-based detector instead of the paper's 10 Hz high-pass, spike removal and Springer HSMM.
\end{document}
